% answers_short.tex —— 2008 年·普通高等学校招生考试·山东卷·文科 · 简短答案速查
% 由 tools/make_short_answers.py 自 answers.tex 的 \gksolhead/\gkanswer 提取，
% 并与 2008/2008_数学真题_山东文_答案要点.md 逐题比对后生成，勿手改（改源 MD 后重跑）。

\begin{gknotice}
\textbf{使用说明}：本节只给最终答案便于快速核对，完整推导见解析版。答案为编写者独立推导，\textbf{非官方标准答案}。
\end{gknotice}

\gkpart{选择题}
\begin{center}\begin{tabular}{ccccccccccccc}
\toprule
题号 & 1 & 2 & 3 & 4 & 5 & 6 & 7 & 8 & 9 & 10 & 11 & 12 \\
\midrule
答案 & B & D & A & C & A & C & D & C & B & C & B & A \\
\bottomrule
\end{tabular}\end{center}

\gkpart{填空题}
\begin{itemize}[leftmargin=2.6em]
  \item[\textbf{13.}] $\dfrac{x^2}4-\dfrac{y^2}{12}=1$。
  \item[\textbf{14.}] $4$。
  \item[\textbf{15.}] $2008$。
  \item[\textbf{16.}] $11$。
\end{itemize}

\gkpart{解答题}
\begin{itemize}[leftmargin=2.6em]
  \item[\textbf{17.}] （1）$\sqrt2$；（2）$\left[\dfrac{\pi}6+k\pi,\dfrac{2\pi}3+k\pi\right]$（$k\in\mathbf Z$）。
  \item[\textbf{18.}] （1）$\dfrac13$；（2）$\dfrac56$。
  \item[\textbf{19.}] （1）证明见解析；（2）$16\sqrt3$。
  \item[\textbf{20.}] （1）$\left\{\dfrac1{S_n}\right\}$ 是公差 $\dfrac12$ 的等差数列，$b_1=1$，$b_n=-\dfrac{2}{n(n+1)}$（$n\geqslant2$）；（2）$-\dfrac{2(2^k-1)}{k(k+1)}$。
  \item[\textbf{21.}] （1）$a=-\dfrac13$，$b=-1$；（2）减区间 $(-\infty,-2)$、$(0,1)$，增区间 $(-2,0)$、$(1,+\infty)$；（3）$f(x)\geqslant g(x)$，等号仅在 $x=0$ 或 $x=1$ 处成立。
  \item[\textbf{22.}] （1）$\dfrac{x^2}5+\dfrac{y^2}4=1$；（2）① $\dfrac{x^2}{4\lambda^2}+\dfrac{y^2}{5\lambda^2}=1$（点 $M$ 不与原点重合）；② $\dfrac{40}9$。
\end{itemize}
